\documentclass[12pt,a4paper]{article}
\usepackage[utf8]{inputenc}
\usepackage[T1]{fontenc}
\usepackage{amsmath}
\usepackage{amssymb}
\usepackage{graphicx}
\usepackage[spanish]{babel}
\usepackage[a4paper, margin=2.5cm]{geometry}
\usepackage[european]{circuitikz}
\usetikzlibrary{babel}

\title{Análisis de Corriente en Circuito Paralelo}
\author{Gemini Code Assist}
\date{\today}

\begin{document}

\maketitle

\section{Descripción del Circuito}

El circuito consta de tres ramas conectadas en paralelo entre un nodo superior, denominado "A", y tierra (GND). El objetivo es determinar la fórmula para calcular la corriente que circula a través de la combinación en serie de las resistencias $R_3$ y $R_{ch}$.

\begin{center}
    \begin{circuitikz}[scale=1.2, transform shape]
        \draw (0,0) node[ground] (gnd) {};
        \draw (0,3) node[circ, label=above:A] (A) {};
        
        % Rama 1: Fuente E y R1
        \draw (A) to[R, l=$R_1$] (-4,3) to[V, v<=$E$] (-4,0) -- (gnd);
        
        % Rama 2: Resistencia R2
        \draw (A) to[R, l=$R_2$] (0,0);
        
        % Rama 3: Resistencias R3 y Rch
        \draw (A) to[R, l=$R_3$] (4,3) to[R, l=$R_{ch}$] (4,0) -- (gnd);
    \end{circuitikz}
\end{center}

\section{Derivación de la Fórmula}

Para resolver este problema, primero calcularemos la tensión en el nodo A ($V_A$) y luego usaremos la Ley de Ohm para encontrar la corriente deseada, que llamaremos $I_{ch}$.

\subsection{Paso 1: Calcular la Tensión en el Nodo A ($V_A$)}

Para encontrar la tensión $V_A$, aplicamos la Ley de Corrientes de Kirchhoff (LCK) en el nodo "A". La ley establece que la suma de las corrientes que salen del nodo es igual a cero.
\begin{equation}
    I_1 + I_2 + I_3 = 0
\end{equation}
Donde $I_1$, $I_2$ e $I_3$ son las corrientes que salen del nodo A por cada una de las tres ramas.
\begin{itemize}
    \item \textbf{Corriente en Rama 1:} $I_1 = \frac{V_A - E}{R_1}$
    \item \textbf{Corriente en Rama 2:} $I_2 = \frac{V_A}{R_2}$
    \item \textbf{Corriente en Rama 3:} $I_3 = \frac{V_A}{R_3 + R_{ch}}$
\end{itemize}

Sustituyendo estas expresiones en la ecuación de la LCK:
\begin{equation}
    \frac{V_A - E}{R_1} + \frac{V_A}{R_2} + \frac{V_A}{R_3 + R_{ch}} = 0
\end{equation}
Agrupando los términos que contienen $V_A$ y despejando, obtenemos la misma expresión que con el Teorema de Millman:
\begin{equation}
    V_A = \frac{\frac{E}{R_1}}{\frac{1}{R_1} + \frac{1}{R_2} + \frac{1}{R_3 + R_{ch}}}
\end{equation}

\subsection{Paso 2: Calcular la Corriente $I_{ch}$}

Una vez conocida la tensión $V_A$, que es la tensión aplicada a la tercera rama, la corriente $I_{ch}$ que circula por ella se calcula directamente con la Ley de Ohm. La resistencia total de esta rama es la suma de las resistencias en serie.

\begin{equation}
    I_{ch} = \frac{V_A}{R_3 + R_{ch}}
\end{equation}

Sustituyendo la expresión de $V_A$ que encontramos en el paso anterior, obtenemos la fórmula final para la corriente $I_{ch}$:
\begin{equation}
    I_{ch} = \frac{1}{R_3 + R_{ch}} \left( \frac{\frac{E}{R_1}}{\frac{1}{R_1} + \frac{1}{R_2} + \frac{1}{R_3 + R_{ch}}} \right)
\end{equation}

\end{document}